\IfFileExists{papers/formalized-vs-literature-preamble.tex}{\documentclass[11pt]{article}
\usepackage[utf8]{inputenc}
\usepackage[T1]{fontenc}
\usepackage{amsmath,amssymb,mathtools}
\usepackage{booktabs,tabularx,array}
\usepackage{enumitem}
\usepackage[margin=1in]{geometry}
\usepackage{xcolor}
\usepackage[hidelinks]{hyperref}
\usepackage{microtype}
\setlength{\parindent}{0pt}
\setlength{\parskip}{0.55em}
\newcommand{\Lean}[1]{\texttt{\nolinkurl{#1}}}
\newcommand{\RepoFile}[1]{\texttt{\nolinkurl{#1}}}
\newcommand{\Class}[1]{\textbf{#1}}
\newcommand{\Top}{\mathord{\top}}
\newcommand{\R}{\mathbb{R}}
\newcommand{\ENN}{\mathbb{R}_{\ge 0}^{\infty}}
\newcommand{\KL}{\operatorname{KL}}
\newcommand{\W}{\mathsf{W}}
\newcommand{\statusline}[2]{%
  \begin{center}\fbox{\begin{minipage}{0.92\linewidth}\textbf{Completion class:} #1\\[2pt]
  \textbf{Strongest caveat:} #2\end{minipage}}\end{center}}
\newenvironment{comparison}{\tabularx{\linewidth}{>{\bfseries}p{0.25\linewidth}X}}{\endtabularx}
\newcommand{\snapshot}{\texttt{902f51aa01a737af68f6feb71c83263cc9b1f4d9}}
}{\documentclass[11pt]{article}
\usepackage[utf8]{inputenc}
\usepackage[T1]{fontenc}
\usepackage{amsmath,amssymb,mathtools}
\usepackage{booktabs,tabularx,array}
\usepackage{enumitem}
\usepackage[margin=1in]{geometry}
\usepackage{xcolor}
\usepackage[hidelinks]{hyperref}
\usepackage{microtype}
\setlength{\parindent}{0pt}
\setlength{\parskip}{0.55em}
\newcommand{\Lean}[1]{\texttt{\nolinkurl{#1}}}
\newcommand{\RepoFile}[1]{\texttt{\nolinkurl{#1}}}
\newcommand{\Class}[1]{\textbf{#1}}
\newcommand{\Top}{\mathord{\top}}
\newcommand{\R}{\mathbb{R}}
\newcommand{\ENN}{\mathbb{R}_{\ge 0}^{\infty}}
\newcommand{\KL}{\operatorname{KL}}
\newcommand{\W}{\mathsf{W}}
\newcommand{\statusline}[2]{%
  \begin{center}\fbox{\begin{minipage}{0.92\linewidth}\textbf{Completion class:} #1\\[2pt]
  \textbf{Strongest caveat:} #2\end{minipage}}\end{center}}
\newenvironment{comparison}{\tabularx{\linewidth}{>{\bfseries}p{0.25\linewidth}X}}{\endtabularx}
\newcommand{\snapshot}{\texttt{902f51aa01a737af68f6feb71c83263cc9b1f4d9}}
}
\title{Mohajerin Esfahani--Kuhn Theorem 4.2: formalized result versus the literature}
\author{}
\date{Formalization snapshot \snapshot}
\begin{document}
\maketitle
\statusline{\Class{A/R}: source-shaped empirical scalar dual with a proved weak direction and an explicit reverse-duality edge}{The Lean theorem uses a support-restricted ambiguity set needed by its total real-valued loss encoding, and it does not derive strong duality from Assumption 4.1.}
\section{Source target}
Theorem~4.2 of Mohajerin Esfahani and Kuhn gives a finite-sample Wasserstein-DRO reformulation whose intermediate scalar dual is
\[
 \inf_{\lambda\ge0}\left\{\lambda\epsilon+
 \frac1N\sum_{i=1}^N\sup_{\xi\in\Xi}
 \bigl(\ell(\xi)-\lambda\lVert\xi-\widehat\xi_i\rVert\bigr)\right\}.
\]
See \RepoFile{prose/distilled_literature/MohajerinEsfahaniKuhn2018_data_driven_wdro.tex}.

\section{Lean declaration}
\Lean{MohajerinEsfahaniKuhn2018.worstCaseExpectation_eq_dual} in
\RepoFile{MohajerinEsfahaniKuhn2018/Basic.lean}.

\section{Formalized content}
Lean proves the weak direction with a support-restricted version of the general OT Lagrangian kernel and proves the empirical integral collapse to a finite average.  The equality is completed from an explicit reverse inequality \texttt{hge}.

The ambiguity set is \Lean{MohajerinEsfahaniKuhn2018.wass1BallΞ}, requiring feasible laws to be supported on $\Xi$.  This is a deliberate fidelity correction: the paper effectively encodes $+\infty$ outside $\Xi$, whereas the Lean loss is a total $\R$-valued function.  Without the support restriction, mass could escape $\Xi$ and violate the displayed dual.

\section{Comparison}
\begin{comparison}
Nominal law & Finite empirical law, faithfully represented.\\
Ground cost & Norm cost, matching the order-one formulation in this declaration.\\
Support & Explicit $Q(\Xi)=1$ restriction in Lean; implicit through the paper's extended loss/domain.\\
Convexity and semicontinuity & Retained as source-facing hypotheses but not used to derive \texttt{hge}.\\
Strong-duality proof & Not formalized from primitive assumptions.\\
Finite convex program & Handled by separate declarations.\\
\end{comparison}

\section{Nonclaims}
This theorem is not yet a complete formalization of Theorem~4.2's derivation.  In particular, it does not prove the needed finite-dimensional convex/LP duality or an extremal representation that would discharge the reverse inequality.
\end{document}
